\documentclass[10pt]{extarticle}
\usepackage[a4paper,margin=1.5cm]{geometry}
\usepackage[T1]{fontenc}
\usepackage{amsmath,amssymb}
\usepackage{graphicx}
\usepackage[export]{adjustbox}
\usepackage{multicol}
\usepackage{enumitem}
\usepackage{booktabs,array,tabularx}
\setlength{\columnsep}{0.7cm}
\setlength{\parindent}{0pt}
\setlength{\parskip}{3pt}
\pagestyle{empty}
\makeatletter
% figure registry (figdefs.tex is generated by the extract step)
\newcommand{\qfigdef}[4]{\@namedef{qf@file@#1}{#2}\@namedef{qf@w@#1}{#3}\@namedef{qf@wide@#1}{#4}}
\newcommand{\qf@check}[1]{\@ifundefined{qf@file@#1}{\PackageError{score2}{Unknown figure id #1}{}}{}}
\newcommand{\qf@img}[4][]{\includegraphics[width=#2\dimexpr\@nameuse{qf@w@#4}pt\relax,max width=#3,max height=0.45\textheight,#1]{\@nameuse{qf@file@#4}}}
% Full-width questions: finalize wraps every question that contains a wide figure in \qfullbegin...\qfullend,
% so the WHOLE question (stem, figures, options) is set across the page in reading order, then the two columns
% resume. The question's head - its number, stem and first figure (\qfullsplit marks the end of the head) -
% is kept on one page, so a number is never left at a page foot with its figure overleaf; the rest may break
% across pages normally, so no big gaps are left. Without a \qfullsplit, the whole question is kept together
% if it is under 45% of a page. Never write these yourself.
\newif\ifqf@full
\newif\ifqf@boxed
\newsavebox{\qf@box}
\newcommand{\qf@open}{\global\qf@fulltrue\par\global\qf@boxedtrue
  \setbox\qf@box\vbox\bgroup\hsize\textwidth\linewidth\textwidth}
\newcommand{\qf@emit}[1]{\par\egroup\global\qf@boxedfalse
  \ifdim\dimexpr\ht\qf@box+\dp\qf@box\relax<#1\textheight\noindent\usebox{\qf@box}\par\else\unvbox\qf@box\fi}
\newcommand{\qfullsplit}{\ifqf@boxed\qf@emit{0.85}\fi}
\newcommand{\qf@close}{\ifqf@boxed\qf@emit{0.45}\fi\global\qf@fullfalse}
\newcommand{\qfullbegin}{\end{multicols}\qf@open}
\newcommand{\qfullend}{\qf@close\begin{multicols}{2}}
\newcommand{\qfullstart}{\qf@open}
\newcommand{\qfulldone}{\qf@close}
% \qfig[scale]{id}: block figure, centred, directly where it is written (size set by the tool).
% Fallback if a wide figure is not inside a full-width question: break out of the columns in place - no floats.
\newcommand{\qfig}[2][1]{\qf@check{#2}%
  \ifqf@full\qf@blk{#1}{#2}\else
  \ifnum\@nameuse{qf@wide@#2}=1 \expandafter\@firstoftwo\else\expandafter\@secondoftwo\fi
  {\end{multicols}\begin{center}\qf@img{#1}{\textwidth}{#2}\end{center}\begin{multicols}{2}}%
  {\qf@blk{#1}{#2}}\fi}
\newcommand{\qf@blk}[2]{\par\smallskip{\centering\qf@img{#1}{\linewidth}{#2}\par}\smallskip}
% \qfigi[scale]{id}: inline figure (inside option lists / tables / minipages); never breaks columns.
\newcommand{\qfigi}[2][1]{\qf@check{#2}\qf@img[valign=c]{#1}{\linewidth}{#2}}
\makeatother
% structure macros - lint reads these to check completeness and figure placement
\newcommand{\papersec}[1]{\par\bigskip{\centering\Large\bfseries PAPER-#1\par}\medskip}
\newcommand{\partsec}[2]{\par\medskip\noindent{\large\bfseries PART-#1: #2}\par\smallskip}
\newcommand{\sect}[1]{\par\smallskip\noindent\textbf{#1}\par}
\newcommand{\qn}[1]{\par\smallskip\noindent\textbf{#1.}~}
\newcommand{\paraq}[2]{\par\smallskip\noindent\textbf{Paragraph for Questions #1 and #2}\par}
\newenvironment{qoptions}{\begin{itemize}[leftmargin=1.8em,itemsep=1pt,topsep=2pt,label={}]}{\end{itemize}}
\newcommand{\opt}[2]{\item[(#1)] #2}
% figure registry (generated by score2_tools.py finalize)
\qfigdef{p1_phy_s1_q01_1}{Score Advanced - Test 04 (04 Jul 2021)-img/p1_phy_s1_q01_1.png}{133.7}{0}
\qfigdef{p1_phy_s1_q02_1}{Score Advanced - Test 04 (04 Jul 2021)-img/p1_phy_s1_q02_1.png}{73.0}{0}
\qfigdef{p1_phy_s1_q05_1}{Score Advanced - Test 04 (04 Jul 2021)-img/p1_phy_s1_q05_1.png}{134.3}{0}
\qfigdef{p1_phy_s1_q06_1}{Score Advanced - Test 04 (04 Jul 2021)-img/p1_phy_s1_q06_1.png}{118.3}{0}
\qfigdef{p1_phy_s1_q07_1}{Score Advanced - Test 04 (04 Jul 2021)-img/p1_phy_s1_q07_1.png}{100.5}{0}
\qfigdef{p1_phy_s1_q07_2}{Score Advanced - Test 04 (04 Jul 2021)-img/p1_phy_s1_q07_2.png}{110.4}{0}
\qfigdef{p1_phy_s1_q07_3}{Score Advanced - Test 04 (04 Jul 2021)-img/p1_phy_s1_q07_3.png}{108.7}{0}
\qfigdef{p1_phy_s1_q07_4}{Score Advanced - Test 04 (04 Jul 2021)-img/p1_phy_s1_q07_4.png}{128.0}{0}
\qfigdef{p1_phy_s1_q08_1}{Score Advanced - Test 04 (04 Jul 2021)-img/p1_phy_s1_q08_1.png}{157.0}{0}
\qfigdef{p1_phy_s1_q10_1}{Score Advanced - Test 04 (04 Jul 2021)-img/p1_phy_s1_q10_1.png}{118.6}{0}
\qfigdef{p1_phy_s2_pre_1}{Score Advanced - Test 04 (04 Jul 2021)-img/p1_phy_s2_pre_1.png}{119.9}{0}
\qfigdef{p1_phy_s2_para13_14_4}{Score Advanced - Test 04 (04 Jul 2021)-img/p1_phy_s2_para13_14_4.png}{136.9}{0}
\qfigdef{p1_phy_s2_q05_1}{Score Advanced - Test 04 (04 Jul 2021)-img/p1_phy_s2_q05_1.png}{104.1}{0}
\qfigdef{p1_phy_s3_q01_2}{Score Advanced - Test 04 (04 Jul 2021)-img/p1_phy_s3_q01_2.png}{79.3}{0}
\qfigdef{p1_phy_s3_q01_3}{Score Advanced - Test 04 (04 Jul 2021)-img/p1_phy_s3_q01_3.png}{70.9}{0}
\qfigdef{p1_phy_s3_q01_4}{Score Advanced - Test 04 (04 Jul 2021)-img/p1_phy_s3_q01_4.png}{71.3}{0}
\qfigdef{p1_phy_s3_q01_5}{Score Advanced - Test 04 (04 Jul 2021)-img/p1_phy_s3_q01_5.png}{52.5}{0}
\qfigdef{p1_phy_s3_q01_6}{Score Advanced - Test 04 (04 Jul 2021)-img/p1_phy_s3_q01_6.png}{66.5}{0}
\qfigdef{p1_phy_s3_q02_1}{Score Advanced - Test 04 (04 Jul 2021)-img/p1_phy_s3_q02_1.png}{93.1}{0}
\qfigdef{p1_phy_s3_q02_B}{Score Advanced - Test 04 (04 Jul 2021)-img/p1_phy_s3_q02_B.png}{129.4}{0}
\qfigdef{p1_phy_s3_q02_C}{Score Advanced - Test 04 (04 Jul 2021)-img/p1_phy_s3_q02_C.png}{100.0}{0}
\qfigdef{p1_phy_s3_q02_D}{Score Advanced - Test 04 (04 Jul 2021)-img/p1_phy_s3_q02_D.png}{98.4}{0}
\qfigdef{p1_chem_s1_q07_1}{Score Advanced - Test 04 (04 Jul 2021)-img/p1_chem_s1_q07_1.png}{111.3}{0}
\qfigdef{p1_chem_s1_q07_2}{Score Advanced - Test 04 (04 Jul 2021)-img/p1_chem_s1_q07_2.png}{108.1}{0}
\qfigdef{p1_chem_s1_q07_3}{Score Advanced - Test 04 (04 Jul 2021)-img/p1_chem_s1_q07_3.png}{107.2}{0}
\qfigdef{p1_chem_s1_q07_4}{Score Advanced - Test 04 (04 Jul 2021)-img/p1_chem_s1_q07_4.png}{141.6}{0}
\qfigdef{p1_chem_s1_q08_1}{Score Advanced - Test 04 (04 Jul 2021)-img/p1_chem_s1_q08_1.png}{148.0}{0}
\qfigdef{p1_chem_s1_q08_2}{Score Advanced - Test 04 (04 Jul 2021)-img/p1_chem_s1_q08_2.png}{106.0}{0}
\qfigdef{p1_chem_s1_q08_3}{Score Advanced - Test 04 (04 Jul 2021)-img/p1_chem_s1_q08_3.png}{105.7}{0}
\qfigdef{p1_chem_s1_q08_4}{Score Advanced - Test 04 (04 Jul 2021)-img/p1_chem_s1_q08_4.png}{96.7}{0}
\qfigdef{p1_chem_s1_q09_1}{Score Advanced - Test 04 (04 Jul 2021)-img/p1_chem_s1_q09_1.png}{101.4}{0}
\qfigdef{p1_chem_s1_q09_A}{Score Advanced - Test 04 (04 Jul 2021)-img/p1_chem_s1_q09_A.png}{35.0}{0}
\qfigdef{p1_chem_s1_q09_B}{Score Advanced - Test 04 (04 Jul 2021)-img/p1_chem_s1_q09_B.png}{42.2}{0}
\qfigdef{p1_chem_s1_q09_C}{Score Advanced - Test 04 (04 Jul 2021)-img/p1_chem_s1_q09_C.png}{37.8}{0}
\qfigdef{p1_chem_s1_q09_4}{Score Advanced - Test 04 (04 Jul 2021)-img/p1_chem_s1_q09_4.png}{46.3}{0}
\qfigdef{p1_chem_s2_pre_1}{Score Advanced - Test 04 (04 Jul 2021)-img/p1_chem_s2_pre_1.png}{101.2}{0}
\qfigdef{p1_chem_s2_para13_14_1}{Score Advanced - Test 04 (04 Jul 2021)-img/p1_chem_s2_para13_14_1.png}{178.8}{0}
\qfigdef{p1_chem_s2_q05_1}{Score Advanced - Test 04 (04 Jul 2021)-img/p1_chem_s2_q05_1.png}{23.2}{0}
\qfigdef{p1_chem_s2_q05_2}{Score Advanced - Test 04 (04 Jul 2021)-img/p1_chem_s2_q05_2.png}{23.2}{0}
\qfigdef{p1_chem_s2_q05_3}{Score Advanced - Test 04 (04 Jul 2021)-img/p1_chem_s2_q05_3.png}{25.2}{0}
\qfigdef{p1_chem_s2_q06_2}{Score Advanced - Test 04 (04 Jul 2021)-img/p1_chem_s2_q06_2.png}{101.7}{0}
\qfigdef{p1_chem_s2_q06_1}{Score Advanced - Test 04 (04 Jul 2021)-img/p1_chem_s2_q06_1.png}{90.5}{0}
\qfigdef{p1_chem_s2_q06_4}{Score Advanced - Test 04 (04 Jul 2021)-img/p1_chem_s2_q06_4.png}{82.0}{0}
\qfigdef{p1_chem_s2_q06_3}{Score Advanced - Test 04 (04 Jul 2021)-img/p1_chem_s2_q06_3.png}{68.9}{0}
\qfigdef{p2_phy_s1_q01_1}{Score Advanced - Test 04 (04 Jul 2021)-img/p2_phy_s1_q01_1.png}{153.5}{0}
\qfigdef{p2_phy_s1_q06_1}{Score Advanced - Test 04 (04 Jul 2021)-img/p2_phy_s1_q06_1.png}{122.1}{0}
\qfigdef{p2_phy_s1_q07_1}{Score Advanced - Test 04 (04 Jul 2021)-img/p2_phy_s1_q07_1.png}{152.8}{0}
\qfigdef{p2_phy_s1_q08_1}{Score Advanced - Test 04 (04 Jul 2021)-img/p2_phy_s1_q08_1.png}{89.3}{0}
\qfigdef{p2_phy_s1_q11_1}{Score Advanced - Test 04 (04 Jul 2021)-img/p2_phy_s1_q11_1.png}{95.0}{0}
\qfigdef{p2_phy_s1_q12_1}{Score Advanced - Test 04 (04 Jul 2021)-img/p2_phy_s1_q12_1.png}{131.2}{0}
\qfigdef{p2_phy_s1_q13_1}{Score Advanced - Test 04 (04 Jul 2021)-img/p2_phy_s1_q13_1.png}{105.2}{0}
\qfigdef{p2_phy_s1_q18_1}{Score Advanced - Test 04 (04 Jul 2021)-img/p2_phy_s1_q18_1.png}{150.6}{0}
\qfigdef{p2_phy_s1_q19_1}{Score Advanced - Test 04 (04 Jul 2021)-img/p2_phy_s1_q19_1.png}{153.8}{0}
\qfigdef{p2_phy_s1_q20_1}{Score Advanced - Test 04 (04 Jul 2021)-img/p2_phy_s1_q20_1.png}{154.7}{0}
\qfigdef{p2_phy_s2_q01_1}{Score Advanced - Test 04 (04 Jul 2021)-img/p2_phy_s2_q01_1.png}{64.6}{0}
\qfigdef{p2_phy_s2_q02_1}{Score Advanced - Test 04 (04 Jul 2021)-img/p2_phy_s2_q02_1.png}{132.6}{0}
\qfigdef{p2_phy_s2_q03_1}{Score Advanced - Test 04 (04 Jul 2021)-img/p2_phy_s2_q03_1.png}{105.2}{0}
\qfigdef{p2_phy_s2_q04_1}{Score Advanced - Test 04 (04 Jul 2021)-img/p2_phy_s2_q04_1.png}{108.1}{0}
\qfigdef{p2_phy_s2_q05_1}{Score Advanced - Test 04 (04 Jul 2021)-img/p2_phy_s2_q05_1.png}{118.1}{0}
\qfigdef{p2_phy_s2_q06_1}{Score Advanced - Test 04 (04 Jul 2021)-img/p2_phy_s2_q06_1.png}{108.6}{0}
\qfigdef{p2_phy_s2_q09_1}{Score Advanced - Test 04 (04 Jul 2021)-img/p2_phy_s2_q09_1.png}{96.8}{0}
\qfigdef{p2_chem_s1_q02_1}{Score Advanced - Test 04 (04 Jul 2021)-img/p2_chem_s1_q02_1.png}{55.4}{0}
\qfigdef{p2_chem_s1_q04_1}{Score Advanced - Test 04 (04 Jul 2021)-img/p2_chem_s1_q04_1.png}{55.1}{0}
\qfigdef{p2_chem_s1_q04_2}{Score Advanced - Test 04 (04 Jul 2021)-img/p2_chem_s1_q04_2.png}{56.0}{0}
\qfigdef{p2_chem_s1_q04_3}{Score Advanced - Test 04 (04 Jul 2021)-img/p2_chem_s1_q04_3.png}{55.8}{0}
\qfigdef{p2_chem_s1_q04_4}{Score Advanced - Test 04 (04 Jul 2021)-img/p2_chem_s1_q04_4.png}{56.0}{0}
\qfigdef{p2_chem_s1_q06_1}{Score Advanced - Test 04 (04 Jul 2021)-img/p2_chem_s1_q06_1.png}{119.3}{0}
\qfigdef{p2_chem_s1_q06_2}{Score Advanced - Test 04 (04 Jul 2021)-img/p2_chem_s1_q06_2.png}{42.4}{0}
\qfigdef{p2_chem_s1_q06_3}{Score Advanced - Test 04 (04 Jul 2021)-img/p2_chem_s1_q06_3.png}{44.2}{0}
\qfigdef{p2_chem_s1_q06_4}{Score Advanced - Test 04 (04 Jul 2021)-img/p2_chem_s1_q06_4.png}{45.4}{0}
\qfigdef{p2_chem_s1_q06_5}{Score Advanced - Test 04 (04 Jul 2021)-img/p2_chem_s1_q06_5.png}{43.7}{0}
\qfigdef{p2_chem_s1_q07_1}{Score Advanced - Test 04 (04 Jul 2021)-img/p2_chem_s1_q07_1.png}{104.0}{0}
\qfigdef{p2_chem_s1_q07_2}{Score Advanced - Test 04 (04 Jul 2021)-img/p2_chem_s1_q07_2.png}{105.2}{0}
\qfigdef{p2_chem_s1_q07_3}{Score Advanced - Test 04 (04 Jul 2021)-img/p2_chem_s1_q07_3.png}{105.2}{0}
\qfigdef{p2_chem_s1_q07_4}{Score Advanced - Test 04 (04 Jul 2021)-img/p2_chem_s1_q07_4.png}{103.7}{0}
\qfigdef{p2_chem_s1_q10_1}{Score Advanced - Test 04 (04 Jul 2021)-img/p2_chem_s1_q10_1.png}{155.9}{0}
\qfigdef{p2_chem_s1_q10_A}{Score Advanced - Test 04 (04 Jul 2021)-img/p2_chem_s1_q10_A.png}{53.7}{0}
\qfigdef{p2_chem_s1_q10_3}{Score Advanced - Test 04 (04 Jul 2021)-img/p2_chem_s1_q10_3.png}{56.0}{0}
\qfigdef{p2_chem_s1_q10_5}{Score Advanced - Test 04 (04 Jul 2021)-img/p2_chem_s1_q10_5.png}{52.8}{0}
\qfigdef{p2_chem_s1_q10_4}{Score Advanced - Test 04 (04 Jul 2021)-img/p2_chem_s1_q10_4.png}{53.0}{0}
\qfigdef{p2_chem_s1_q11_1}{Score Advanced - Test 04 (04 Jul 2021)-img/p2_chem_s1_q11_1.png}{96.5}{0}
\qfigdef{p2_chem_s1_q11_2}{Score Advanced - Test 04 (04 Jul 2021)-img/p2_chem_s1_q11_2.png}{98.3}{0}
\qfigdef{p2_chem_s1_q11_3}{Score Advanced - Test 04 (04 Jul 2021)-img/p2_chem_s1_q11_3.png}{97.2}{0}
\qfigdef{p2_chem_s1_q11_4}{Score Advanced - Test 04 (04 Jul 2021)-img/p2_chem_s1_q11_4.png}{97.9}{0}
\qfigdef{p2_chem_s1_q15_1}{Score Advanced - Test 04 (04 Jul 2021)-img/p2_chem_s1_q15_1.png}{134.0}{0}
\qfigdef{p2_chem_s1_q15_A}{Score Advanced - Test 04 (04 Jul 2021)-img/p2_chem_s1_q15_A.png}{116.9}{0}
\qfigdef{p2_chem_s1_q15_B}{Score Advanced - Test 04 (04 Jul 2021)-img/p2_chem_s1_q15_B.png}{118.1}{0}
\qfigdef{p2_chem_s1_q15_C}{Score Advanced - Test 04 (04 Jul 2021)-img/p2_chem_s1_q15_C.png}{118.6}{0}
\qfigdef{p2_chem_s1_q15_5}{Score Advanced - Test 04 (04 Jul 2021)-img/p2_chem_s1_q15_5.png}{119.3}{0}
\qfigdef{p2_chem_s1_q16_1}{Score Advanced - Test 04 (04 Jul 2021)-img/p2_chem_s1_q16_1.png}{140.6}{0}
\qfigdef{p2_chem_s1_q16_B}{Score Advanced - Test 04 (04 Jul 2021)-img/p2_chem_s1_q16_B.png}{141.1}{0}
\qfigdef{p2_chem_s1_q16_3}{Score Advanced - Test 04 (04 Jul 2021)-img/p2_chem_s1_q16_3.png}{139.5}{0}
\qfigdef{p2_chem_s1_q16_4}{Score Advanced - Test 04 (04 Jul 2021)-img/p2_chem_s1_q16_4.png}{141.6}{0}
\qfigdef{p2_chem_s1_q18_1}{Score Advanced - Test 04 (04 Jul 2021)-img/p2_chem_s1_q18_1.png}{155.8}{0}
\qfigdef{p2_chem_s1_q19_1}{Score Advanced - Test 04 (04 Jul 2021)-img/p2_chem_s1_q19_1.png}{140.6}{0}
\qfigdef{p2_chem_s1_q19_2}{Score Advanced - Test 04 (04 Jul 2021)-img/p2_chem_s1_q19_2.png}{138.8}{0}
\qfigdef{p2_chem_s1_q19_3}{Score Advanced - Test 04 (04 Jul 2021)-img/p2_chem_s1_q19_3.png}{137.3}{0}
\qfigdef{p2_chem_s1_q19_4}{Score Advanced - Test 04 (04 Jul 2021)-img/p2_chem_s1_q19_4.png}{118.8}{0}
\qfigdef{p2_chem_s1_q20_A}{Score Advanced - Test 04 (04 Jul 2021)-img/p2_chem_s1_q20_A.png}{30.6}{0}
\qfigdef{p2_chem_s2_q02_1}{Score Advanced - Test 04 (04 Jul 2021)-img/p2_chem_s2_q02_1.png}{65.3}{0}
\qfigdef{p2_chem_s2_q04_1}{Score Advanced - Test 04 (04 Jul 2021)-img/p2_chem_s2_q04_1.png}{155.5}{0}
\raggedcolumns
\makeatletter\@beginparpenalty=10000\makeatother
\newcolumntype{L}[1]{>{\raggedright\arraybackslash}p{#1}}
\begin{document}
\begin{center}{\Large\bfseries Score Advanced --- Test 04 (04 Jul 2021)}\end{center}
\begin{multicols}{2}
\raggedcolumns
\papersec{1}
\partsec{1}{PHYSICS}
\sect{SEC-1}
\qn{1} Figure shows small part of circuit. [Given C$_1$ = C$_2$ = C$_3$ = C$_4$ = 5 $\mu$F] Which of the following statement is/are true ?\par
\qfig{p1_phy_s1_q01_1}
\begin{qoptions}
\opt{A}{Potential of point O can be determined but that of B can't be determined.}
\opt{B}{Potential of both point O \& B can't be determined.}
\opt{C}{If charge on capacitor C$_2$ were also specified, potential of point O can be determined but that of B cann't be determined.}
\opt{D}{If charge on capacitor C$_2$ were also specified, potential of both the point O \& B can be determined.}
\end{qoptions}
\qn{2} A cylinder having length 10cm, and density 0.65 g/cm$^3$ floats half submerged in the liquid shown. ($\rho_{\text{air}}$ = 1.3 $\times$ 10$^{-3}$ g/cm$^3$) Now pressure of air made 100 time of initial pressure keeping temperature constant. Which of the following statement(s) is/are true :\par
\qfig{p1_phy_s1_q02_1}
\begin{qoptions}
\opt{A}{The cylinder moves upwards.}
\opt{B}{The cylinder moves downwards.}
\opt{C}{Displacement of the cylinder is 0.56 cm.}
\opt{D}{Displacement of the cylinder is 0.60 cm.}
\end{qoptions}
\qn{3} Assume that the earth changes its shape and turns into an infinite cylinder whose radius and distance of moon from the axis of cylindrical earth remains unchanged and also moon is remain spherical. Assume density of earth remains same and distance of moon from earth is much greater than radius of earth. Which of the following statement(s) is/are correct :
\begin{qoptions}
\opt{A}{Gravitational field at the moon due to earth increases.}
\opt{B}{Gravitational field at the moon due to earth decreases.}
\opt{C}{Speed of the moon in its orbit around the cylindrical earth increases.}
\opt{D}{Speed of the moon in its orbit around the cylindrical earth decreases.}
\end{qoptions}
\qn{4} A particle of mass m is thrown in a parabolic path towards earth having mass M and radius R. It comes to a greatest approach distance R/2 from surface of earth. Magnitude of total energy change required to make the particle moving in an elliptical path having eccentricity 1/2, with the greatest approach point either aphelion or perihelion are :
\begin{qoptions}
\opt{A}{$\dfrac{1}{3}\dfrac{\text{GMm}}{R}$}
\opt{B}{$\dfrac{1}{6}\dfrac{\text{GMm}}{R}$}
\opt{C}{$\dfrac{1}{2}\dfrac{\text{GMm}}{R}$}
\opt{D}{$\dfrac{\text{GMm}}{R}$}
\end{qoptions}
\qn{5} An air-filled parallel plate capacitor with the plate area A is connected to a battery with an emf V volt and negligible internal resistance. One of the plates vibrates so that the distance between plates varies as d = d$_0$ + a cos$\omega$t (where a $\ll$ d$_0$). The capacitor breaks down when the instantaneous current reaches to the value of I. Choose the correct options(s).\par
\qfig{p1_phy_s1_q05_1}
\begin{qoptions}
\opt{A}{The maximum possible amplitude vibrations `a' is $\dfrac{\text{Id}_0^2}{\text{VA}\omega\varepsilon_0}$}
\opt{B}{The maximum possible amplitude vibrations `a' is $\dfrac{2\text{Id}_0^2}{\text{VA}\omega\varepsilon_0}$}
\opt{C}{The instantaneous current in circuit is $\dfrac{\varepsilon_0\text{V a}}{(\text{d}_0+\text{a cos}\omega\text{t})^{\,2}}\sin\omega t$.}
\opt{D}{The instantaneous current in circuit is $\dfrac{\varepsilon_0\text{V a}}{2(\text{d}_0+\text{a cos}\omega\text{t})^{\,2}}\sin\omega t$.}
\end{qoptions}
\qn{6} A small solid ball of density $\rho$ is held inside at point A of cubical container of side L, filled with an ideal liquid of density 4$\rho$ as shown in the figure. Now, if the container starts moving with constant acceleration a horizontally and the ball is released from point A simultaneously, then\par
\qfig{p1_phy_s1_q06_1}
\begin{qoptions}
\opt{A}{For ball to hit the top of container at end Q, a = 3g}
\opt{B}{For ball to hit the top of container at end Q, a = 2g}
\opt{C}{Ball hits the top of container at end Q after a time t = $\sqrt{\dfrac{L}{3g}}$}
\opt{D}{Ball hits the top of container at end Q after a time t = $\sqrt{\dfrac{2L}{3g}}$}
\end{qoptions}
\qn{7} A time varying voltage is applied across AB as shown. Potential difference across the capacitor plate varies as shown. Then which of the following is/are correct.\par
\qfig{p1_phy_s1_q07_1}
\qfig{p1_phy_s1_q07_2}
\begin{qoptions}
\opt{A}{Plot of potential difference across CD is \qfigi{p1_phy_s1_q07_3}}
\opt{B}{Plot of potential difference across CD is \qfigi{p1_phy_s1_q07_4}}
\opt{C}{Plot of current across resistance does not vary for large value of resistance \& capacitor}
\opt{D}{Plot of current across resistance does not vary for small value of resistance \& capacitor}
\end{qoptions}
\qn{8} A cavity is taken out from a uniform conducting sphere. Inside the cavity a dipole is placed as shown in the figure. Find which of the following is/are true (Given : q = 10$^{-1}$C, $\ell$ = 0.1 mm, $\theta$ = 30$^\circ$, d = 10 cm, R = 12 cm)\par
\qfig{p1_phy_s1_q08_1}
\begin{qoptions}
\opt{A}{Potential at points A \& B are same}
\opt{B}{Potential at points A \& B are different}
\opt{C}{Potential at point P is 0 volt}
\opt{D}{Potential at point P is 2 volt}
\end{qoptions}
\qn{9} Water jet coming out of a stationary horizontal tube at speed v = 2m/s strikes horizontally a massive wall moving in opposite direction with same speed. Water comes to rest relative to wall after striking. Cross-section of jet is 3 cm$^2$ and density of water is 1 gm/cc. Choose the correct statement(s) :-
\begin{qoptions}
\opt{A}{Force exerted on the wall is 1.2 N}
\opt{B}{Force exerted on the wall is 2.4 N}
\opt{C}{Rate of change of kinetic energy of water jet striking the wall is 19.2 J/s.}
\opt{D}{Rate of change of kinetic energy of water jet striking the wall is zero.}
\end{qoptions}
\qn{10} A potentiometer as shown in figure in which three batteries between A \& B are connected when the point P is the null point, measured by the potentiometer arrangement. Then choose correct statement(s)\par
\qfig{p1_phy_s1_q10_1}
\begin{qoptions}
\opt{A}{Current passing through 6 V battery is 4A}
\opt{B}{Current passing through 2 V battery is 4A}
\opt{C}{V$_\text{A}$ $-$ V$_\text{B}$ = 2V}
\opt{D}{V$_\text{A}$ $-$ V$_\text{B}$ = 6V}
\end{qoptions}
\sect{SEC-2}
\paraq{11}{12}
When a conductor is charged then its potential rises. The increase in potential is directly proportional to the charge given to the conductor.\par
Q $\propto$ V $\Rightarrow$ Q = CV\par
The constant C is known as the capacity of the conductor.\par
Let C is capacitance of a conductor. On being connected to a battery. It charges to a potential V from zero potential. If q is charge on the conductor at that time then q = CV. Let battery supplies small amount of charge dq to the conductor at constant potential V. Then small amount of work done by the battery against the force exerted by exsiting charge is\par
$\text{dW} = \text{V dq} = \dfrac{q}{C}\text{dq} \Rightarrow \text{W} = \displaystyle\int_0^Q \dfrac{q}{C}\,\text{dq} = \dfrac{1}{C}\left[\dfrac{q^2}{2}\right]_0^Q \Rightarrow \text{W} = \dfrac{Q^2}{2C}$\par
Two capacitors of capacitance C$_1$ = 2$\mu$F \& C$_2$ = 4$\mu$F are charged to Potential difference 100 V \& 150V respectively. The charged capacitors are connected as shown in figure through a battery \& resistance. At time t = 0 switch S is closed.\par
\qfig{p1_phy_s2_pre_1}
\qn{1} Choose the correct statement(s) :
\begin{qoptions}
\opt{A}{Current through the battery immediately after closing switch is 10 A.}
\opt{B}{Current through the battery immediately after closing switch is 3 A.}
\opt{C}{work done by battery is $\dfrac{20}{3}$ mJ}
\opt{D}{work done by battery is $\dfrac{10}{3}$ mJ}
\end{qoptions}
\qn{2} Final charges on each capacitor will be :
\begin{qoptions}
\opt{A}{$\dfrac{2000}{3}\,\mu$C, $\dfrac{1400}{3}\,\mu$C}
\opt{B}{$\dfrac{1000}{3}\,\mu$C, $\dfrac{1800}{3}\,\mu$C}
\opt{C}{$\dfrac{1200}{3}\,\mu$C, $\dfrac{800}{3}\,\mu$C}
\opt{D}{$\dfrac{1000}{3}\,\mu$C, $\dfrac{1400}{3}\,\mu$C}
\end{qoptions}
\paraq{13}{14}
A light body revolving around a heavier body due to gravitational attraction, is called satellite. Earth is a satellite of the Sun while Moon is satellite of Earth.\par
\textbf{Orbital velocity (v$_0$) :} A satellite of mass m moving in an orbit of radius r with speed v$_0$ then required centripetal force is provided by gravitation.\par
$\text{F}_{\text{cp}} = \text{F}_{\text{g}} \Rightarrow \dfrac{mv_0^2}{r} = \dfrac{\text{GMm}}{r^2} \Rightarrow v_0 = \sqrt{\dfrac{\text{GM}}{r}} = \sqrt{\dfrac{\text{GM}}{(R_e+h)}}$, (r = R$_e$ + h)\par
For a satellite very close to the Earth surface h $\ll$ R$_e$ $\therefore$ r = R$_e$\par
$v_0 = \sqrt{\dfrac{\text{GM}}{R_e}} = \sqrt{gR_e} = 8$ km/s\par
Essential Condition's for Satellite Motion\par
Centre of satellite's orbit coincide with centre of Earth.\par
Plane of orbit of satellite is passing through centre of Earth.\par
\qfig{p1_phy_s2_para13_14_4}
\qn{3} A near earth orbit satellite of earth experiences very small air resistance. It's period of revolution decreases by 1.08 sec/day.
\begin{qoptions}
\opt{A}{It's radius decreases by $\dfrac{160}{3}$ m in 1 day.}
\opt{B}{It's radius decreases by $\dfrac{160}{3}$ m in 1 revolution.}
\opt{C}{It's angular momentum decreases}
\opt{D}{It's mechanical energy increases.}
\end{qoptions}
\qn{4} Satellite A is in a circular orbit of radius r and satellite B in another circular orbit of radius 4 r, both revolving around the earth. The masses of the satellites A and B are in the ratio 3 : 1. If the symbols V, E, u and T represent the speed, total energy, escape velocity and period of revolution of satellite, then
\begin{qoptions}
\opt{A}{V$_\text{A}$/V$_\text{B}$ = 2}
\opt{B}{E$_\text{A}$/E$_\text{B}$ = 12}
\opt{C}{u$_\text{A}$/u$_\text{B}$ = 2}
\opt{D}{T$_\text{A}$/T$_\text{B}$ = $\dfrac{1}{8}$}
\end{qoptions}
\paraq{15}{16}
Electric potential is a scalar property of every point in the region of electric field. At a point in electric field potential is defined as the interaction energy of a unit positive charge. If at a point in electric field a charge q$_0$ has potential energy U, then electric potential at that point can be given as\par
$\text{V} = \dfrac{U}{q_0}$ joule/coulomb\par
Potential difference between two points in electric field can be defined as work done in displacing a unit positive charge from one point to another against the electric forces.\par
If a unit +ve charge is displaced from a point A to B as shown work required can be given as\par
$\text{V}_\text{B} - \text{VA} = -\displaystyle\int_A^B \vec{E}\cdot d\vec{r}$
\qn{5} Consider three equipotential surfaces of field produced by $-$q point charge. A small point charge +q is taken slowly along two paths from A to B and A to C, then :-\par
\qfig{p1_phy_s2_q05_1}
\begin{qoptions}
\opt{A}{Work done by electric field in moving it from A to B is +ve}
\opt{B}{Work done by electric field in moving it from A to B is $-$ve}
\opt{C}{Work done by an external agent in moving it from A to C is +ve}
\opt{D}{The correct order of potential will be V$_3$ $>$ V$_2$ $>$ V$_1$.}
\end{qoptions}
\qn{6} The potential of electric field V = x$^2$y, where x and y denotes coordinate of a point in electric field region. Which of the following is incorrect.
\begin{qoptions}
\opt{A}{Electric field at all points on positive x-axis is in positive y-direction}
\opt{B}{Electric field at all points on positive y-axis is in negative x-direction.}
\opt{C}{Electric field at origin is non-zero.}
\opt{D}{Electric field at (x, y) makes an angle of $\theta = \cot^{-1}\left(\dfrac{2y}{x}\right)$ with x-axis.}
\end{qoptions}
\sect{SEC-3}
\qn{1} \noindent{\footnotesize\setlength{\tabcolsep}{2pt}\begin{tabular}{|L{.5cm}|L{2.0cm}|L{.6cm}|L{3.6cm}|}\hline
 & Column-I & & Column-II\\\hline
(A) & Magnitude of angular momentum of system about O increases after the instant shown. & (P) & A dipole is hinged on a long smooth fixed horizontal wire. A charge Q (A) is threaded on wire and is free to move on the wire. The system comprises of the dipole and point charge. Gravity is absent.\par\qfigi{p1_phy_s3_q01_2}\\\hline
(B) & The net torque on body A is zero about O at the instant shown. & (Q) & A system of two charges $-$Q and 2Q (A) of mass m each are connected by a light rod and kept on smooth level ground.An external uniform electric field E$_0$ = mg/Q is applied as shown. Gravity is present.\par\qfigi{p1_phy_s3_q01_3}\\\hline
(C) & The net force on negative charge increases in magnitude after the instant shown. & (R) & An electron (A) is revolving about a fixed proton in anticlockwise direction The system comprises of the electron and proton. Gravity is absent.\par\qfigi{p1_phy_s3_q01_4}\\\hline
(D) & Mechanical energy of the system increase after the instant shown. (ignore potential energy of external electrical field.) & (S) & System consists of a particle (A) of mass m and charge $-$Q. It is projected from level ground in presence of constant electric field E = mg/Q in vertically upward direction. Gravity is present.\par\qfigi{p1_phy_s3_q01_5}\\\hline
 & & (T) & System consists of a uniform ring and particle (A) of mass m and charge $-$Q on its circumference as shown. The ring is in pure rolling on a sufficiently rough horizontal surface. Gravity is present.\par\qfigi{p1_phy_s3_q01_6}\\\hline
\end{tabular}}\par
\qn{2} In column-I, we have certain situations where a battery of variable voltage V is connected in a circuit. Column-II gives comments on current in element A on variation of voltage :-\par
\noindent{\footnotesize\setlength{\tabcolsep}{2pt}\begin{tabular}{|L{.5cm}|L{3.4cm}|L{.6cm}|L{2.4cm}|}\hline
 & Column-I & & Column-II\\\hline
(A) & \qfigi{p1_phy_s3_q02_1} & (P) & on increasing V, i may/will increase\\\hline
(B) & \qfigi{p1_phy_s3_q02_B} & (Q) & On increasing V, i may/will decrease\\\hline
(C) & \qfigi{p1_phy_s3_q02_C} & (R) & On increasing V, i may/will not change at all\\\hline
(D) & \qfigi{p1_phy_s3_q02_D} & (S) & On reversing V, i may/will be zero\\\hline
 & & (T) & On reversing V, i may/will increase\\\hline
\end{tabular}}\par
\partsec{2}{CHEMISTRY}
\sect{SEC-1}
\qn{1} 10 gm of an impure sample of CaCO$_3$ on heating gives 4.48 gm of solid residue. Assuming that all CaCO$_3$ is completely decomposed into CaO and CO$_2$, identify correct statement(s) \{Ca = 40\}
\begin{qoptions}
\opt{A}{The impurity is non-volatile}
\opt{B}{Impurity is volatile}
\opt{C}{Percent purity of CaCO$_3$ sample is 80\%}
\opt{D}{CO$_2$ liberated is 3.52 gm}
\end{qoptions}
\qn{2} The following four reactions are at equilibrium in closed rigid containers. In which of the reaction(s) equilibrium is not disturbed by adding He gas.
\begin{qoptions}
\opt{A}{H$_2$(g) + I$_2$(g) $\rightleftharpoons$ 2HI(g)}
\opt{B}{N$_2$O$_4$(g) $\rightleftharpoons$ 2NO$_2$(g)}
\opt{C}{CO(g) + 2H$_2$(g) $\rightleftharpoons$ CH$_3$OH(g)}
\opt{D}{C(s) + H$_2$O(g) $\rightleftharpoons$ CO(g) + H$_2$(g)}
\end{qoptions}
\qn{3} A and B are two volatile liquids with their vapour pressure : P$^0_\text{A}$ = 340mm of Hg and P$^0_\text{B}$ = 420mm of Hg at 25$^\circ$C. To a flask containing 8 mole of `A', 5 mole of `B' was added. As soon as `B' was added, `A' starts associating to form a non-volatile solid A$_\text{x}$(solid) which is soluble in both liquids A and B. The association of `A' follows 1$^{\text{st}}$ order kinetics. The vapour pressure of solution measured after 1 hour is 360 mm of Hg and after very long time it is 300 mm of Hg. Identify correct statements(s)
\begin{qoptions}
\opt{A}{The value of x = 4}
\opt{B}{At t = 1 hr, 0.5 mole of A$_\text{x}$ (solid) is formed}
\opt{C}{At t = 1 hr, 4 mole of `A' (liquid) remains in solution.}
\opt{D}{The value of x = 2}
\end{qoptions}
\qn{4} Which of the following order(s) is/are correct?
\begin{qoptions}
\opt{A}{Al $>$ Ga (Atomic Size)}
\opt{B}{P$^{-3}$ $>$ Ca$^{+2}$ (Ionic Size)}
\opt{C}{Ne $>$ Al (Electron Affinity)}
\opt{D}{Li $>$ B (First Ionisation Energy)}
\end{qoptions}
\qn{5} Select the reaction(s) which is/are not possible.
\begin{qoptions}
\opt{A}{2NaF + 2H$_2$O $\rightarrow$ 2NaOH + H$_2$ + F$_2$}
\opt{B}{KHF$_2$ (fused) $\xrightarrow{\text{electrolysis}}$ H$_2$ + F$_2$}
\opt{C}{NaCl + H$_2$O $\rightarrow$ NaOH + HCl}
\opt{D}{2NaOH (hot and conc.) + Cl$_2$ $\longrightarrow$ NaCl + NaOCl + H$_2$O}
\end{qoptions}
\qn{6} Which of the following element(s) undergoes disproportination in basic medium?
\begin{qoptions}
\opt{A}{P$_4$}
\opt{B}{Cl$_2$}
\opt{C}{S$_8$}
\opt{D}{Si}
\end{qoptions}
\qn{7} Which of the following pair of compounds doesn't represent enantiomeric pair :
\begin{qoptions}
\opt{A}{\qfigi{p1_chem_s1_q07_1}}
\opt{B}{\qfigi{p1_chem_s1_q07_2}}
\opt{C}{\qfigi{p1_chem_s1_q07_3}}
\opt{D}{\qfigi{p1_chem_s1_q07_4}}
\end{qoptions}
\qn{8} Identify pair(s) of compounds which can be seperate by fractional distillation ?
\begin{qoptions}
\opt{A}{\qfigi{p1_chem_s1_q08_1}}
\opt{B}{\qfigi{p1_chem_s1_q08_2}}
\opt{C}{\qfigi{p1_chem_s1_q08_3}}
\opt{D}{\qfigi{p1_chem_s1_q08_4}}
\end{qoptions}
\qn{9}\par
\qfig{p1_chem_s1_q09_1}
Possible product(s) is/are :
\begin{qoptions}
\opt{A}{\qfigi{p1_chem_s1_q09_A}}
\opt{B}{\qfigi{p1_chem_s1_q09_B}}
\opt{C}{\qfigi{p1_chem_s1_q09_C}}
\opt{D}{\qfigi{p1_chem_s1_q09_4}}
\end{qoptions}
\qn{10} Correct statement :
\begin{qoptions}
\opt{A}{Vinyl chloride $>$ Benzyl chloride \quad (C$-$Cl bond length)}
\opt{B}{Allyl bromide $>$ Vinyl bromide \quad (Reactivity towards S$_\text{N}$2)}
\opt{C}{Tertiary butyl bromide $>$ Secondary butyl bromide \quad (Reactivity towards S$_\text{N}$1)}
\opt{D}{7-Chloro 1,3,5-cycloheptatriene $>$ 5-chloro 1,3-cyclopentadiene \quad (Rate of hydrolysis)}
\end{qoptions}
\sect{SEC-2}
\paraq{11}{12}
Titration of acid samples with hydroxide ion solution produces very different pH curves depending upon the type and strength of acid. Given below are two titration curves for two different monobasic acid having same initial concentration of 0.1M\par
\qfig{p1_chem_s2_pre_1}
\qn{1} With respect to above information in paragraph indentify correct statement(s)
\begin{qoptions}
\opt{A}{Curve (1) may be of HCl while (2) may be of CH$_3$COOH.}
\opt{B}{The nature of solution at equivalence point for curve(1) is acidic}
\opt{C}{Phenolphthalein can be used to detect the equivalence point for curve (1)}
\opt{D}{For weaker acid among (1) and (2) the value of K$_\text{a}$ = 10$^{-5}$}
\end{qoptions}
\qn{2} Let the colour of an acidic indicator HIn (colourless) be visible only when it's ionized form (pink) is 25\% or more than the unionized form in a solution. If HIn (pK$_\text{a}$ = 9) is added to solution of pH = 9.6, then correct observation(s) is/are -
\begin{qoptions}
\opt{A}{Pink colour will be visible}
\opt{B}{Pink colour will not be visible}
\opt{C}{Unionized form of indicator is less than 25\% in solution}
\opt{D}{Ionized form of indicator is less than 25\% in solution}
\end{qoptions}
\paraq{13}{14}
\qfig{p1_chem_s2_para13_14_1}
\qn{3} The gas C can be differentiate from SO$_2$ by using of :-
\begin{qoptions}
\opt{A}{KMnO$_4$/H$^+$}
\opt{B}{K$_2$Cr$_2$O$_7$/H$^+$}
\opt{C}{KIO$_3$ + starch}
\opt{D}{KOH(aq.)}
\end{qoptions}
\qn{4} Which of the following is/are correct for E(s)?
\begin{qoptions}
\opt{A}{It reacts with D to form a very hard substance used as an abrasive}
\opt{B}{Most abundant element in earth crust}
\opt{C}{It reacts with oxygen at 1173 K to form a 3D silicate}
\opt{D}{It can be purified by zone refining}
\end{qoptions}
\paraq{15}{16}
Compound A$_1$, A$_2$, A$_3$ having molecular formula C$_6$H$_{10}$. All have five membered ring. A$_1$, A$_2$, A$_3$ are strutural isomers.\par
A$_1$ $\xrightarrow[\text{(CCl}_4\text{)}]{\text{Br}_2}$ Enantiomeric pair is formed as product\par
A$_2$ $\xrightarrow[\text{(CCl}_4\text{)}]{\text{Br}_2}$ Diastereomeric pair is formed as product\par
A$_3$ $\xrightarrow[\text{(CCl}_4\text{)}]{\text{Br}_2}$ Non-resolvable optically inactive product is formed
\qn{5} A$_1$ can be :
\begin{qoptions}
\opt{A}{\qfigi{p1_chem_s2_q05_1}}
\opt{B}{\qfigi{p1_chem_s2_q05_2}}
\opt{C}{\qfigi{p1_chem_s2_q05_3}}
\opt{D}{All}
\end{qoptions}
\qn{6} A$_2$ \& A$_3$ may respectively :
\begin{qoptions}
\opt{A}{\qfigi{p1_chem_s2_q06_2}}
\opt{B}{\qfigi{p1_chem_s2_q06_1}}
\opt{C}{\qfigi{p1_chem_s2_q06_4}}
\opt{D}{\qfigi{p1_chem_s2_q06_3}}
\end{qoptions}
\sect{SEC-3}
\qn{1} Match the column :\par
\noindent{\footnotesize\setlength{\tabcolsep}{2pt}\begin{tabular}{|L{.5cm}|L{2.6cm}|L{.6cm}|L{3.2cm}|}\hline
 & Column I & & Column II\\\hline
(A) & CH$_3$COCH$_3$ + CHCl$_3$ & (P) & $\Delta$S$_\text{mix}$ $>$ 0\\\hline
(B) & C$_2$H$_5$OH + H$_2$O & (Q) & $\Delta$V$_\text{mix}$ $>$ 0\\\hline
(C) & C$_2$H$_5$Br + C$_2$H$_5$I & (R) & $\Delta$H$_\text{mix}$ $<$ 0\\\hline
(D) & CH$_3$COCH$_3$ + CS$_2$ & (S) & Maximum boiling azeotrope\\\hline
 & & (T) & Minimum boiling azeotrope\\\hline
\end{tabular}}\par
\qn{2} Match the Column.\par
\noindent{\footnotesize\setlength{\tabcolsep}{2pt}\begin{tabular}{|L{.5cm}|L{2.4cm}|L{.6cm}|L{3.2cm}|}\hline
 & Column-I Compounds & & Column-II Characteristics\\\hline
(A) & Tetrachloroethene & (P) & Liquefaction of it is due to London forces only\\\hline
(B) & NO$_2$ & (Q) & Used as a solvent for dry cleaning\\\hline
(C) & NO & (R) & At least one atom is sp$^2$ hybridized\\\hline
(D) & CO & (S) & Bond order of at least one bond is greater than 2\\\hline
 & & (T) & Odd electron species\\\hline
\end{tabular}}\par
\partsec{3}{MATHEMATICS}
\sect{SEC-1}
\qn{1} If the pair of lines $x^2+2xy+ay^2=0$ and $ax^2+2xy+y^2=0$ have exactly one line in common and the joint equation of the other two lines is $\alpha x^2+\beta xy+\gamma y^2=0$ where $\alpha,\beta,\gamma\in N$, then lowest value of
\begin{qoptions}
\opt{A}{$\alpha=3$}
\opt{B}{$\beta=9$}
\opt{C}{$\alpha+\gamma=6$}
\opt{D}{$\alpha+\beta=13$}
\end{qoptions}
\qn{2} Two vertices Q and R of a $\Delta$PQR are mirror images of each other with respect to the line y = x. Q is (5,4) and the orthocentre of $\Delta$PQR is origin. If P is (a,b) then [a + b] is greater than\par
[Note : [ ] denotes greatest integer function]
\begin{qoptions}
\opt{A}{6}
\opt{B}{7}
\opt{C}{8}
\opt{D}{9}
\end{qoptions}
\qn{3} The equation of a plane which is equally inclined to the lines $\dfrac{x-1}{2}=\dfrac{y}{-2}=\dfrac{z+2}{-1}$ and $\dfrac{x+3}{8}=\dfrac{y-4}{1}=\dfrac{z}{-4}$ and passing origin is/are
\begin{qoptions}
\opt{A}{$14x-5y-7z=0$}
\opt{B}{$3x-4y-z=0$}
\opt{C}{$x+2y-5z=0$}
\opt{D}{$2x-7y-z=0$}
\end{qoptions}
\qn{4} Let the intercept made by the line $\dfrac{x}{1}=\dfrac{y}{2}=\dfrac{z}{3}$ between the plane $x+y+z=3$ and the line $\dfrac{x-1}{2}=\dfrac{y-2}{-1}=\dfrac{z-3}{-1}$ is k, then
\begin{qoptions}
\opt{A}{$\log_{14}(2k)=\dfrac{1}{3}$}
\opt{B}{$\log_{14}(2k)=\dfrac{1}{2}$}
\opt{C}{$\log_2(2k)\in(2,3)$}
\opt{D}{$\log_2(2k)\in(3,4)$}
\end{qoptions}
\qn{5} Let $\vec{p}=x^2\hat{i}-3\hat{j}+(x+3)\hat{k}$ and $\vec{q}=\hat{i}+3\hat{j}-(x-3)\hat{k}$ be two vectors such that $|\vec{p}|=|\vec{q}|$. If $\vec{r}=2\vec{p}+3\vec{q}$ and $\vec{s}=3\vec{p}-2\vec{q}$, then
\begin{qoptions}
\opt{A}{$\displaystyle\int_1^3(\vec{r}\,{}^{\wedge}\vec{s})\,dx=\dfrac{\pi}{2}$}
\opt{B}{$\displaystyle\int_1^3(\vec{r}\,{}^{\wedge}\vec{s})\,dx=\pi$}
\opt{C}{$\displaystyle\int_2^{5/2}(\vec{r}\,{}^{\wedge}\vec{s})\,dx=\dfrac{\pi}{2}$}
\opt{D}{$\displaystyle\int_2^{5/2}(\vec{r}\,{}^{\wedge}\vec{s})\,dx=\dfrac{\pi}{4}$}
\end{qoptions}
\qn{6} If the curves $x^2-y^2=4$ and $xy=\sqrt{5}$ intersects at points A and B, then the possible number of point(s) C on curve $x^2-y^2=4$ such that $\Delta$ABC is equilateral is less than or equal to-
\begin{qoptions}
\opt{A}{0}
\opt{B}{1}
\opt{C}{2}
\opt{D}{4}
\end{qoptions}
\qn{7} If number of integral points ($\alpha,\beta$) (an integral point is a point whose both co-ordinates are integers) for which the area common to the circle $x^2+y^2+2\alpha^2x+2\beta^2y+c=0$ and its image in the line $x+y+4=0$ is maximum is k, then k is greater than or equal to-
\begin{qoptions}
\opt{A}{3}
\opt{B}{4}
\opt{C}{5}
\opt{D}{6}
\end{qoptions}
\qn{8} 2 parabolas have the focus (5,$-$2). Their directrices are x-axis and y-axis respectively. If the slope of their common chord is m then m$^2$ is less
\begin{qoptions}
\opt{A}{5}
\opt{B}{3}
\opt{C}{$\sqrt{2}$}
\opt{D}{1}
\end{qoptions}
\qn{9} If the planes $\vec{r}\cdot(\hat{i}+\hat{j}+\hat{k})=q_1,\ \vec{r}\cdot(\hat{i}+2a\hat{j}+\hat{k})=q_2$ and $\vec{r}\cdot(a\hat{i}+a^2\hat{j}+\hat{k})=q_3$ intersect in a line then the value of `a' is(are)
\begin{qoptions}
\opt{A}{$\dfrac{1}{4}$}
\opt{B}{$\dfrac{1}{3}$}
\opt{C}{$\dfrac{1}{2}$}
\opt{D}{1}
\end{qoptions}
\qn{10} Let $\vec{a}$ and $\vec{b}$ are two vectors such that $\vec{a}=2(\sin\alpha+\cos\alpha)\hat{i}+\hat{j}$ and $\vec{b}=(\sin\beta)\hat{i}+(\cos\beta)\hat{j}$ where $\alpha$ and $\beta$ satisfy the relation $2(\sin\alpha+\cos\alpha)\sin\beta=3-\cos\beta$, then\par
(where $\alpha,\beta\in R$ and $\alpha\in\left[0,\dfrac{\pi}{2}\right)$)
\begin{qoptions}
\opt{A}{$2\tan^2\alpha+\tan^2\beta=9$}
\opt{B}{$\tan^2\alpha+2\tan^2\beta=9$}
\opt{C}{$\tan^2\alpha+4\sqrt{2}\tan\beta=17$}
\opt{D}{$\tan^2\alpha-\tan^2\beta=1$}
\end{qoptions}
\sect{SEC-2}
\paraq{11}{12}
The circle $C_1: x^2+y^2=r^2$ and ellipse $C_2: \dfrac{x^2}{16}+\dfrac{y^2}{9}=1$ intersect at 4 distinct points A,B,C,D. Their common tangents form a parallelogram A$'$B$'$C$'$D$'$.
\qn{1} If A$'$B$'$C$'$D$'$ is a square then r is equal to $\dfrac{a}{\sqrt{b}}$ where a and b are coprime then -
\begin{qoptions}
\opt{A}{a + b = 12}
\opt{B}{a + b = 7}
\opt{C}{a + 2b = 9}
\opt{D}{a + b$^2$ = 9}
\end{qoptions}
\qn{2} If A$'$B$'$C$'$D$'$ is a square, then the ratio of area of the circle C$_1$ to the area of the circumcircle of $\Delta$A$'$B$'$C$'$ is K, then K is greater than or equal to
\begin{qoptions}
\opt{A}{$\dfrac{9}{16}$}
\opt{B}{$\dfrac{3}{4}$}
\opt{C}{$\dfrac{1}{2}$}
\opt{D}{$\dfrac{1}{4}$}
\end{qoptions}
\paraq{13}{14}
Let A,B,C be 3 sets such that\par
$A=\left\{(x,y):\dfrac{x}{\cos\theta}=\dfrac{y}{\sin\theta}=2,\ \theta\in R\right\}$\par
$B=\left\{(x,y):\dfrac{x+1}{\cos2205^\circ}=\dfrac{y}{\sin2205^\circ}=r,\ r\in R\right\}$\par
$C=\left\{(x,y):(x-\sqrt{3})^2+(y-1)^2\leqslant a^2;\ a\in R\right\}$\par
Let C$_1$,C$_2$ be the curves represented by A and B respectively while R be the region represented by C.
\qn{3} A points P(x$_1$,y$_1$) and Q(x$_2$,y$_2$) be such that P,Q $\in$ A $\cap$ B, then point of intersection of tangents drawn to C$_1$ at P and Q is-
\begin{qoptions}
\opt{A}{(2,3)}
\opt{B}{($-$4,4)}
\opt{C}{(3,$-$3)}
\opt{D}{($-$5,5)}
\end{qoptions}
\qn{4} If minimum value of `a' for which A $\cap$ C = A is $\alpha$, then length of curve C$_2$ lying in R corresponding to a = $\alpha$, is
\begin{qoptions}
\opt{A}{$\sqrt{50}$}
\opt{B}{$\sqrt{52}$}
\opt{C}{$\sqrt{58}$}
\opt{D}{$\sqrt{57}$}
\end{qoptions}
\paraq{15}{16}
Let P$_1$ : x + y + 2z $-$ 3 = 0 and P$_2$ : x $-$ 2y + z = 4 be two planes. Let A(1,3,4) and B(3,2,7) be two points in R$^3$.
\qn{5} The equation of plane P$_3$ through line of intersection of P$_1$ and P$_2$ and length of projection upon which of the line segment AB is greatest is equal to ax + by + cz + 1 = 0, then (where a,b,c $\in$ I)
\begin{qoptions}
\opt{A}{a + b + c = 4}
\opt{B}{a + 2b + c = 7}
\opt{C}{a $-$ b + c = 0}
\opt{D}{a + 2b + c = 5}
\end{qoptions}
\qn{6} The equation of plane P$_4$ through the line of intersection of P$_1$ and P$_2$ and the length of projection upon which the line segment AB is least is
\begin{qoptions}
\opt{A}{2x $-$ y + 3z $-$ 7 = 0}
\opt{B}{3y $-$ z + 1 = 0}
\opt{C}{x + 3y + z + 2 = 0}
\opt{D}{3x $-$ 3y + 4z $-$ 11 = 0}
\end{qoptions}
\sect{SEC-3}
\qn{1} Let two lines be L$_1$ : $\dfrac{x-1}{2}=\dfrac{y-2}{3}=\dfrac{z-\alpha}{4}$ and L$_2$ : $\dfrac{x-3}{\beta}=\dfrac{y-5}{6}=\dfrac{z-5}{8}$, then\par
\noindent{\footnotesize\setlength{\tabcolsep}{2pt}\begin{tabular}{|L{.5cm}|L{2.8cm}|L{.6cm}|L{3.0cm}|}\hline
 & Column-I & & Column-II\\\hline
(A) & If $\beta$ = 4 and $\alpha\in R-\{1\}$ & (P) & Non coplanar lines\\\hline
(B) & If $\beta$ = 4 and $\alpha$ = 1 & (Q) & Coplanar lines\\\hline
(C) & If $\beta\in R-\{4\}$ and $\alpha$ = 1 & (R) & Lines lie in a unique plane\\\hline
(D) & If $\beta\in R-\{4\}$ and $\alpha\in R-\{1\}$ & (S) & Infinite planes containing both the lines\\\hline
 & & (T) & Lines do not intersect\\\hline
\end{tabular}}\par
\qn{2} \noindent{\footnotesize\setlength{\tabcolsep}{2pt}\begin{tabular}{|L{.5cm}|L{4.2cm}|L{.6cm}|L{1.2cm}|}\hline
 & Column-I & & Column-II\\\hline
(A) & If a variable line y = ax + b is tangent to ellipse 3x$^2$ + 4y$^2$ = 6, then locus of point P(a,b) is a conic whose eccentricity is equal to & (P) & $\dfrac{\sqrt{3}}{2}$\\\hline
(B) & Consider three vectors $\vec{a},\vec{b}$ and $\vec{c}$ such that $|\vec{a}|=\sqrt{3},\vec{a}\cdot\vec{c}=2$ and $\vec{a}\times\vec{c}=\vec{b}$, then value of $\dfrac{1}{8}\left|\vec{a}\times\vec{b}+3\vec{c}+2\vec{a}\right|$ is equal to & (Q) & $\sqrt{\dfrac{3}{2}}$\\\hline
(C) & Let A and B denote the points of intersection of the circles x$^2$ + y$^2$ $-$ 6x + 4y = 3 and x$^2$ + y$^2$ + 4x $-$ 4y = 17, then slope of line AB is equal to & (R) & 1\\\hline
(D) & Locus of the point of intersection of the pair of perpendicular tangents to the circles x$^2$ + y$^2$ = 1 and x$^2$ + y$^2$ = 7 is a circle of radius `r' then `r' is greater than & (S) & 2\\\hline
 & & (T) & $\dfrac{5}{4}$\\\hline
\end{tabular}}\par
\raggedcolumns
\papersec{2}
\partsec{1}{PHYSICS}
\sect{SEC-A}
\qn{1} The output of given logic circuit is :-\par
\qfig{p2_phy_s1_q01_1}
\begin{qoptions}
\opt{A}{A $\cdot$ (B + C)}
\opt{B}{A$\cdot$ (B$\cdot$ C)}
\opt{C}{(A + B)$\cdot$ (A + C)}
\opt{D}{A + B + C}
\end{qoptions}
\qn{2} Which one of the following statements is wrong?
\begin{qoptions}
\opt{A}{Radio waves in the frequency range 10MHz to 20 MHz are called sky waves}
\opt{B}{Radio horizon of the transmitting antenna for space wave is d$_\text{T}$ = $\sqrt{(\text{Rh}_\text{T})}$ (R = radius of earth, h$_\text{T}$ = height of transmitting antenna)}
\opt{C}{Fiber optical communication is almost free from electrical disturbances}
\opt{D}{The principle of fibre optical communication is total internal reflection}
\end{qoptions}
\qn{3} Three copper rods are subjected to different potential difference. Compare the drift speed of electrons through them. Assume that all 3 are at the same temperature.\par
\noindent{\footnotesize\begin{tabular}{|l|l|l|l|}\hline
 & Length & Diameter & Potential difference\\\hline
(A) & L & 3d & V\\\hline
(B) & 2L & d & 2V\\\hline
(C) & 3L & 2d & 2V\\\hline
\end{tabular}}
\begin{qoptions}
\opt{A}{v$_\text{A}$ = v$_\text{B}$ $>$ v$_\text{C}$}
\opt{B}{v$_\text{A}$ $>$ v$_\text{B}$ $>$ v$_\text{C}$}
\opt{C}{v$_\text{A}$ $<$ v$_\text{B}$ $<$ v$_\text{C}$}
\opt{D}{None of these}
\end{qoptions}
\qn{4} Three infinitely long linear charges of charge density $\lambda$, $\lambda$ and $-2\,\lambda$ are placed in space. A point in space is specified by its perpendicular distance r$_1$, r$_2$ \& r$_3$ respectively from the linear charges. For the points which are equipotential :-
\begin{qoptions}
\opt{A}{$\dfrac{\text{r}_1\text{r}_2}{\text{r}_3^2}$ = constant}
\opt{B}{r$_1$r$_2$r$_3^2$ = constant}
\opt{C}{r$_1$r$_2$r$_3^{1/2}$ = constant}
\opt{D}{r$_1$r$_2$r$_3$ = constant}
\end{qoptions}
\qn{5} Choose the correct option :-
\begin{qoptions}
\opt{A}{Work done by external agent to take negative charge from higher potential to lower potential is negative.}
\opt{B}{Work done by external agent to take positive charge from lower potential to higher potential is positive.}
\opt{C}{Work done by electric field on negative charge for motion from higher potential energy to lower potential energy is negative.}
\opt{D}{Work done by electric field on positive charge for motion from lower potential energy to higher potential energy is positive.}
\end{qoptions}
\qn{6} Figure shows a siphon. Choose the wrong statement: (P$_0$ = atmospheric pressure)\par
\qfig{p2_phy_s1_q06_1}
\begin{qoptions}
\opt{A}{Siphon works when h$_3$ $>$ 0}
\opt{B}{Pressure at point 2 is P$_2$ = P$_0$ $-$ $\rho$gh$_3$}
\opt{C}{Pressure at point 3 is P$_0$}
\opt{D}{Pressure at point 1 is P$_0$}
\end{qoptions}
\qn{7} Rank in order, from highest to lowest, the liquid heights h$_\text{a}$ to h$_\text{d}$. The air flow is from left to right. The liquid columns are not drawn to scale :-\par
\qfig{p2_phy_s1_q07_1}
\begin{qoptions}
\opt{A}{h$_\text{a}$ $>$ h$_\text{b}$ = h$_\text{d}$ $>$ h$_\text{c}$}
\opt{B}{h$_\text{b}$ = h$_\text{d}$ $>$ h$_\text{c}$ $>$ h$_\text{a}$}
\opt{C}{h$_\text{d}$ $>$ h$_\text{c}$ $>$ h$_\text{b}$ $>$ h$_\text{a}$}
\opt{D}{h$_\text{b}$ $>$ h$_\text{d}$ $>$ h$_\text{c}$ $>$ h$_\text{a}$}
\end{qoptions}
\qn{8} A cylinder of radius 4cm and height 10cm is immersed in two liquids as shown. Specific gravity of oil is 0.5. 2cm of cylinder is in the air. Select the incorrect statement. Neglect atmospheric pressure.\par
\qfig{p2_phy_s1_q08_1}
\begin{qoptions}
\opt{A}{Force exerted by oil on the cylinder is zero.}
\opt{B}{Force exerted by water on the cylinder is 0.96 $\pi$ newton.}
\opt{C}{Force exerted by both water and oil is 1.14 $\pi$ newton.}
\opt{D}{Frequency of oscillation does not depend on density of oil, if cylinder is displaced slightly.}
\end{qoptions}
\qn{9} A 2 $\mu$F capacitor is charged to a potential = 10V. Another 4 $\mu$F capacitor is charged to a potential = 20V. The two capacitors are then connected in a single loop, with the positive plate of one connected with negative plate of the other. What heat is evolved in the circuit ?
\begin{qoptions}
\opt{A}{300 $\mu$J}
\opt{B}{600 $\mu$J}
\opt{C}{900 $\mu$J}
\opt{D}{450 $\mu$J}
\end{qoptions}
\qn{10} Time period of simple pendulum increases by an amount $\sqrt{2}$ times at height `h' from the surface of earth. Then the value of h is :
\begin{qoptions}
\opt{A}{0.25 R}
\opt{B}{0.4 R}
\opt{C}{0.5 R}
\opt{D}{R}
\end{qoptions}
\qn{11} Reverse break down voltage for silicon diode is 6V for a particular doping concentration and forward break down voltage is 0.3V then find approximate current through battery in the given circuit\par
\qfig{p2_phy_s1_q11_1}
\begin{qoptions}
\opt{A}{2 Amp}
\opt{B}{1.3 Amp}
\opt{C}{1 Amp}
\opt{D}{$\dfrac{\pi}{\sqrt{2}}$ Amp}
\end{qoptions}
\qn{12} A particle of mass 1gm and charge $-$0.1$\mu$C is projected from ground with a velocity 10$\sqrt{2}$ at an 45$^\circ$ with horizontal in the area having uniform electric field 1 kV/cm in horizontal direction. Acceleration due to gravity is 10m/s$^2$ in vertical downward direction. Select incorrect statement:\par
\qfig{p2_phy_s1_q12_1}
\begin{qoptions}
\opt{A}{Time of flight of particle is 2 sec}
\opt{B}{Range of particle is 20 m}
\opt{C}{Total displacement of particle is 0m}
\opt{D}{Particle will follow straight line motion.}
\end{qoptions}
\qn{13} Two light bulbs in the circuit have rating A [24 V, 48W] and B [24V, 36W] as shown.\par
\qfig{p2_phy_s1_q13_1}
\begin{qoptions}
\opt{A}{When switch is closed bulb A delivers 24 watt \& B delivers 18 watt.}
\opt{B}{When switch is open A delivers 9 watt \& B delivers 12 watt approximately.}
\opt{C}{When switch is open A delivers 24 watt \& B delivers 18 watt approximately.}
\opt{D}{When switch is open `A' delivers 6 watt and B delivers 9 watt approximately.}
\end{qoptions}
\qn{14} A 0.5 kg mass of lead is submerged in a container filled to the brim with water and a block of wood floats on top. The lead mass is slowly lifted from the container by a thin wire and as it emerges into air the level of the water in the container drops a bit. The lead mass is now placed on the block of wood. As the lead is placed on the wood.
\begin{qoptions}
\opt{A}{Some water spills over the edge of the container.}
\opt{B}{The water level rises exactly to the brim as before.}
\opt{C}{The water level rises but does not reach the brim of the container}
\opt{D}{There is not enough information provided to decide between the three options.}
\end{qoptions}
\qn{15} A spherical asteroid having the same density as that of earth is floating in free space. A small pebble is revolving around the asteroid under the influence of gravity near the surface of the asteroid. What is the approximate time period of the pebble?
\begin{qoptions}
\opt{A}{24 h}
\opt{B}{365 days}
\opt{C}{10 min}
\opt{D}{1 hr 24 min}
\end{qoptions}
\qn{16} A transistor is connected in common emitter (CE) configuration. The collector supply is 8 V and the voltage drop across a resistor of 800 $\Omega$ in the collector circuit is 0.5 V. If the current gain factor $\alpha$ is 0.96, find the base current.
\begin{qoptions}
\opt{A}{48 $\mu$A}
\opt{B}{44 $\mu$A}
\opt{C}{26 $\mu$A}
\opt{D}{15 $\mu$A}
\end{qoptions}
\qn{17} Carbon resistor has resistance specified by three bands having colour red, yellow and black. If this resistor is cut into two pieces of equal length then the new colour code of each one will be (Neglect tolerance of 4th band)
\begin{qoptions}
\opt{A}{Brown, Red, Black}
\opt{B}{Red, Orange, Black}
\opt{C}{White, Blue, Black}
\opt{D}{Black, Blue, Black}
\end{qoptions}
\qn{18} Two resistors with resistances R$_1$ and R$_2$ are connected in series and so are two capacitors with capacitances C$_1$ and C$_2$, these two systems are connected to a battery as shown in figure. Potential difference between a and b is zero in steady state if\par
\qfig{p2_phy_s1_q18_1}
\begin{qoptions}
\opt{A}{R$_1$ + R$_2$ = $\dfrac{\text{C}_1\text{C}_2}{\text{C}_1+\text{C}_2}$}
\opt{B}{R$_1$R$_2$ = C$_1$C$_2$}
\opt{C}{R$_1$C$_2$ = R$_2$C$_1$}
\opt{D}{R$_1$C$_1$ = R$_2$C$_2$}
\end{qoptions}
\qn{19} A mass m, travelling at speed V$_0$ in a straight line from far away is deflected when it passes near a black hole of mass M which is at a perpendicular distance R from the original line of flight. a, the distance of closest approach between the mass and the black hole is given by the relation :\par
\qfig{p2_phy_s1_q19_1}
\begin{qoptions}
\opt{A}{$a=R\left(1+\dfrac{2GM}{aV_0^2}\right)^{1/2}$}
\opt{B}{$a=R\left(1+\dfrac{aV_0^2}{2GM}\right)^{1/2}$}
\opt{C}{$a=R\left(1+\dfrac{GM}{2aV_0^2}\right)^{-1/2}$}
\opt{D}{$a=R\left(1+\dfrac{2GM}{aV_0^2}\right)^{-1/2}$}
\end{qoptions}
\qn{20} Circuit given below is that of\par
\qfig{p2_phy_s1_q20_1}
\begin{qoptions}
\opt{A}{AND gate}
\opt{B}{OR gate}
\opt{C}{NAND gate}
\opt{D}{NOR gate}
\end{qoptions}
\sect{SEC-B}
\qn{1} The figure shows a hollow hemisphere of radius R in which two charges 3q and 5q are placed symmetrically about the centre O on the planar surface. The electric flux over the curved surface is $\dfrac{nq}{\varepsilon_0}$. Fill the value of n.\par
\qfig{p2_phy_s2_q01_1}
\qn{2} Circuit for the measurement of resistance by potentiometer is shown. The galvanometer is first connected at point A and zero deflection observed at length P J = 10 cm. In second case it is connected at point C and zero deflection is observed at a length 30cm from P. Then the unknown resistance X is nR, fill the value of n.\par
\qfig{p2_phy_s2_q02_1}
\qn{3} Four ammeters with identical internal resistances r and a resistor of resistance R are connected to a current source as shown in figure. It is known that the reading of the ammeter A$_1$ is I$_1$ = 3 A and the reading of the ammeter A$_2$ is I$_2$ = 5 A. Determine the ratio of the resistances R/r.\par
\qfig{p2_phy_s2_q03_1}
\qn{4} The value of current through the 20$\Omega$ resistor is $\dfrac{n}{10}$ amp. Fill the value of n.\par
\qfig{p2_phy_s2_q04_1}
\qn{5} During lighting, a current pulse, shown in figure, flows from the cloud at a height 1.5km to the ground. If the breakdown electric field of humid air is about 400 kVm$^{-1}$, the energy released during lighting would be (in unit of 10$^9$ J)\par
\qfig{p2_phy_s2_q05_1}
\qn{6} A hemispherical bowl of mass m is uniformly charged with charge density `$\sigma$'. Electric potential due to charge distribution at a point `A' is $\dfrac{\sigma R}{n\varepsilon_0}$, fill the value of n. (which lies at centre of radius as shown).\par
\qfig{p2_phy_s2_q06_1}
\qn{7} Dipole having dipole moment 4C-m is placed on origin along x-axis. Point charge of 8$\mu$C is fixed at (4, 0, 0) m. Now dipole is rotated by an angle of $\dfrac{\pi}{2}$. If the work done (in kJ) in rotating the dipole is 6x, then fill the value of x.
\qn{8} A soap bubble in a form of circular tube having radius of curvature R and radius of curvature perpendicular to it is 5R. If the excess pressure in the bubble is $\dfrac{3\text{nT}}{5R}$, then fill the value of n.
\qn{9} A 1$\mu$F capacitor is connected in the circuit shown below. The emf of the cell is 3 volts and internal resistance is 0.5 ohms. The resistors R$_1$ and R$_2$ have values 4 ohms and 1 ohm respectively. The charge (in $\mu$C) on the capacitor in steady state must be :\par
\qfig{p2_phy_s2_q09_1}
\qn{10} A parallel plate capacitor of capacitance C is connected to a battery and is charged to a potential difference V. Another capacitor of capacitance 2C is similarly charged to a potential difference 2V. The charging battery is now disconnected and the capacitors are now connect in parallel to each other in such a way that the positive terminal of one is connected to the negative terminal of the other. The final energy of the configuration is $\dfrac{\text{nCV}^2}{2}$, then fill the value of n.
\partsec{2}{CHEMISTRY}
\sect{SEC-A}
\qn{1} Observe the following four solutions :\par
P : 10 ml, 0.1 M NaOH + 5ml, 0.1 M HCl\par
Q : 10 ml, 0.1 M NaOH + 15ml, 0.1 M CH$_3$COOH\par
R : 10 ml, 0.1 M NH$_3$ + 10ml, 0.1 M NH$_4$Cl\par
S : 10 ml, 0.05 M NaF + 5ml, 0.1 M HF\par
Which of above solutions act as buffer
\begin{qoptions}
\opt{A}{P and Q only}
\opt{B}{R and S only}
\opt{C}{P, R, S only}
\opt{D}{Q, R, S only}
\end{qoptions}
\qn{2} Which of the following air pollutant is likely not to produce photo chemical smog?
\begin{qoptions}
\opt{A}{NO}
\opt{B}{CH$_2$CHCHO}
\opt{C}{\qfigi{p2_chem_s1_q02_1}}
\opt{D}{CO$_2$}
\end{qoptions}
\qn{3} Which of the following is false regarding acid rain.
\begin{qoptions}
\opt{A}{Acid rain has pH below 5.6}
\opt{B}{CO$_2$ and CO are major contributors of acid rain}
\opt{C}{Acid rain washes away plant nutrients needed for their growth}
\opt{D}{Acid rain is responsible for the change of colour of Tajmahal.}
\end{qoptions}
\qn{4} An equilibrium shift towards reactants at higher temperatures. Find the correct graph :
\begin{qoptions}
\opt{A}{\qfigi{p2_chem_s1_q04_1}}
\opt{B}{\qfigi{p2_chem_s1_q04_2}}
\opt{C}{\qfigi{p2_chem_s1_q04_3}}
\opt{D}{\qfigi{p2_chem_s1_q04_4}}
\end{qoptions}
\qn{5} Which of the following properties is/are true for hydrogen peroxide?\par
(i) Mild disinfectant\par
(ii) Rocket propellant\par
(iii) Bleaching agent\par
(iv) Oxidising agent
\begin{qoptions}
\opt{A}{only (iii), (iv)}
\opt{B}{only (i), (iii) (iv)}
\opt{C}{only (i), (iv)}
\opt{D}{All (i), (ii), (iii) \& (iv)}
\end{qoptions}
\qn{6} Major product of reaction is :\par
\qfig{p2_chem_s1_q06_1}
\begin{qoptions}
\opt{A}{\qfigi{p2_chem_s1_q06_2}\par + it's enantiomer}
\opt{B}{\qfigi{p2_chem_s1_q06_3}}
\opt{C}{\qfigi{p2_chem_s1_q06_4}}
\opt{D}{\qfigi{p2_chem_s1_q06_5}}
\end{qoptions}
\qn{7} Which of following reaction sequence shows overall inversion of configuration in product ?
\begin{qoptions}
\opt{A}{\qfigi{p2_chem_s1_q07_1}}
\opt{B}{\qfigi{p2_chem_s1_q07_2}}
\opt{C}{\qfigi{p2_chem_s1_q07_3}}
\opt{D}{\qfigi{p2_chem_s1_q07_4}}
\end{qoptions}
\qn{8} The reaction for which K$_\text{C}$ $\leq$ K$_\text{p}$ at 298 K
\begin{qoptions}
\opt{A}{3O$_2$(g) $\rightleftharpoons$ 2O$_3$(g)}
\opt{B}{N$_2$(g) + 3H$_2$(g) $\rightleftharpoons$ 2NH$_3$(g)}
\opt{C}{N$_2$(g) + 2NO$_2$(g) $\rightleftharpoons$ 2NO$_2$(g)}
\opt{D}{H$_2$(g) + I$_2$(g) $\rightleftharpoons$ 2HI(g)}
\end{qoptions}
\qn{9} Which of the following order is NOT correct, against mentioned properties?
\begin{qoptions}
\opt{A}{Sc $>$ Y $>$ La (Ionisation energy)}
\opt{B}{Hf$^{2+}$ $<$ Hf $\simeq$ Zr (size)}
\opt{C}{A$\ell$ $>$ Si (Electronegativity)}
\opt{D}{Ga $>$ A$\ell$ (Ionisation energy)}
\end{qoptions}
\qn{10} Major product (Q) of following reaction is :\par
\qfig{p2_chem_s1_q10_1}
\begin{qoptions}
\opt{A}{\qfigi{p2_chem_s1_q10_A}}
\opt{B}{\qfigi{p2_chem_s1_q10_3}}
\opt{C}{\qfigi{p2_chem_s1_q10_5}}
\opt{D}{\qfigi{p2_chem_s1_q10_4}}
\end{qoptions}
\qn{11} Which of following reaction(s) produce Saytzeff product as a major product :
\begin{qoptions}
\opt{A}{\qfigi{p2_chem_s1_q11_1}}
\opt{B}{\qfigi{p2_chem_s1_q11_2}}
\opt{C}{\qfigi{p2_chem_s1_q11_3}}
\opt{D}{\qfigi{p2_chem_s1_q11_4}}
\end{qoptions}
\qn{12} Which of the following liquid pair does not show contraction in volume on mixing -
\begin{qoptions}
\opt{A}{CHCl$_3$+CH$_3$COCH$_3$}
\opt{B}{H$_2$O + HCl}
\opt{C}{H$_2$O + HNO$_3$}
\opt{D}{H$_2$O + C$_2$H$_5$OH}
\end{qoptions}
\qn{13} In which of the following process energy is absorbed?
\begin{qoptions}
\opt{A}{S$^{2-}_{(g)}$ $\rightarrow$ S$^{-}_{(g)}$ + e$^-$}
\opt{B}{N$^{-}_{(g)}$ $\rightarrow$ N$_{(g)}$ + e$^-$}
\opt{C}{O$^{-}_{(g)}$ + e$^-$ $\rightarrow$ O$^{2-}_{(g)}$}
\opt{D}{Na$^{+}_{(g)}$ + e$^-$ $\rightarrow$ Na$_{(g)}$}
\end{qoptions}
\qn{14} Consider the following reactions :-\par
(i) F$_2$ + 2X$-\!\!-\!\!\rightarrow$ 2F$^-$ + X$_2$ (X = Cl, Br, I)\par
(ii) Cl$_2$ + 2X$^-\!\!-\!\!\rightarrow$ 2Cl$^-$ + X$_2$ (X = Br, I)\par
(iii) I$_2$ + 2Br$^-\!\!-\!\!\rightarrow$ Br$_2$ + 2I$^-$\par
option which have only the correct set of reaction is.
\begin{qoptions}
\opt{A}{(ii), (iii)}
\opt{B}{(i), (iii)}
\opt{C}{(i), (ii), (iii)}
\opt{D}{(i), (ii)}
\end{qoptions}
\qn{15}\par
\qfig{p2_chem_s1_q15_1}
In above reaction out of the following which will be obtained in maximum yield.
\begin{qoptions}
\opt{A}{\qfigi{p2_chem_s1_q15_A}}
\opt{B}{\qfigi{p2_chem_s1_q15_B}}
\opt{C}{\qfigi{p2_chem_s1_q15_C}}
\opt{D}{\qfigi{p2_chem_s1_q15_5}}
\end{qoptions}
\qn{16} Select incorrect option :
\begin{qoptions}
\opt{A}{\qfigi{p2_chem_s1_q16_1}}
\opt{B}{\qfigi{p2_chem_s1_q16_B}}
\opt{C}{\qfigi{p2_chem_s1_q16_3}}
\opt{D}{\qfigi{p2_chem_s1_q16_4}}
\end{qoptions}
\qn{17} Which of the following statement is INCORRECT.
\begin{qoptions}
\opt{A}{In Haber's process iron oxide, potassium oxide and aluminium oxide are used as catalyst.}
\opt{B}{ClO$_2$ is used as bleaching agent}
\opt{C}{Decomposition of ozone is an endothermic process}
\opt{D}{Hydrolysis of XeF$_6$ is a non redox process}
\end{qoptions}
\qn{18}\par
\qfig{p2_chem_s1_q18_1}
Select best reagent to carry out above conversion.
\begin{qoptions}
\opt{A}{LiAlH$_4$}
\opt{B}{NaBH$_4$}
\opt{C}{NH$_2$$-$NH$_2$+H$_2$O$_2$}
\opt{D}{H$_2$ (excess) / Ni, $\Delta$}
\end{qoptions}
\qn{19} Select option in which incorrect major product is shown.
\begin{qoptions}
\opt{A}{\qfigi{p2_chem_s1_q19_1}}
\opt{B}{\qfigi{p2_chem_s1_q19_2}}
\opt{C}{\qfigi{p2_chem_s1_q19_3}}
\opt{D}{\qfigi{p2_chem_s1_q19_4}}
\end{qoptions}
\qn{20} Which of the following decolorises Br$_2$ water solution :
\begin{qoptions}
\opt{A}{\qfigi{p2_chem_s1_q20_A}}
\opt{B}{CH$_3$$-$CH$_2$$-$OH}
\opt{C}{HC $\equiv$ CH}
\opt{D}{CH$_3$$-$CH$_2$$-$Cl}
\end{qoptions}
\sect{SEC-B}
\qn{1} Phosphoric acid is prepared commercially from phosphate ore as following :\par
Ca$_3$(PO4)$_2$ + SiO$_2$ + C + O$_2$ + H$_2$O $-\!\!\rightarrow$ CaSiO$_3$ + CO$_2$ + H$_3$PO$_4$\par
If 1.0 kg each of calcium phosphate and silica are used with excess of C, O$_2$ and H$_2$O, what is the maximum quantity of phosphoric acid (in kg) that can be produced. (Ca = 40, P = 31, Si = 28)
\qn{2} How many 1,2-shift(s) take(s) place in the following carbocation.\par
\qfig{p2_chem_s2_q02_1}
\qn{3} The solubility of silver chromate is 1.992 $\times$ 10$^{-2}$ gm /litre. The molar mass of Ag$_2$CrO$_4$ is 332 gm/mol. If the molar concentration of Ag$^+$ in saturated solution of Ag$_2$CrO$_4$ is X $\times$ 10$^{-4}$ M, give `X'.
\qn{4} Consider the following flow diagram :-\par
\qfig{p2_chem_s2_q04_1}
Give the total sum of Boron atoms present in X and Y respectively.
\qn{5} A fixed amount of glucose is dissolved in 100 gm water to form a solution that freezes at $-$ 0.2$^\circ$C. The solution is cooled down to $-$0.25$^\circ$C. Find how many grams of ice would have separated.
\qn{6} Give the sum of group number and period of element ``Uub''.
\qn{7} The Henry's law constant for solubility of N$_2$ gas in water at 25$^\circ$C is 1 $\times$ 10$^5$ atm. The mole fraction of N$_2$ in air is 0.8. The number of moles of N$_2$ from air that will dissolve in 10 moles of H$_2$O at 25$^\circ$C and at 5 atm pressure is X $\times$ 10$^{-4}$. Give `X'-
\qn{8} 200 ml of a mixture of 50\% H$_2$, 40\% CH$_4$ and 10\% CO$_2$ is exploded with excess of oxygen in an eudiometer tube. The resultant mixture is cooled to room temperature and is passed through aq. KOH. Find volume contraction observed (in ml) after passing through KOH.
\qn{9} 100 ml, 1M H$_2$SO$_4$ solution of density 1.5 gm/ml is mixed with 400 ml of water , resulting in a diluted solution of density 1.25 gm/ml. Find molarity of final diluted solution. (Density of water = 1 gm/ml).
\qn{10} If the number of S$-$S, S$-$O$-$S and O$-$H linkages in a molecule of pyrosulphurous acid are x, y and z respectively.\par
Give (x + y + z).
\partsec{3}{MATHEMATICS}
\sect{SEC-A}
\qn{1} Set of values of `$\alpha$' for which the angle bisector of the lines $(\alpha+1)x+2y+5=0$ and $4x+\alpha y-3=0$ containing origin is also the obtuse angle bisector, is-
\begin{qoptions}
\opt{A}{$\left(-\infty,-\dfrac{2}{3}\right)$}
\opt{B}{$\left(-\dfrac{2}{3},\infty\right)$}
\opt{C}{$\left(-\infty,-\dfrac{2}{3}\right)\cup(1,\infty)$}
\opt{D}{$(-1,\infty)$}
\end{qoptions}
\qn{2} Let mirror image of point A(1,4) in line y = x is point B; mirror image of point B in line y = $-$x is C and mirror image of C in x-axis is D, then area of triangle ABD is-
\begin{qoptions}
\opt{A}{3 sq. units}
\opt{B}{$\dfrac{3}{2}$ sq. units}
\opt{C}{6 sq. units}
\opt{D}{0 sq. units}
\end{qoptions}
\qn{3} Let P be a variable point on the ellipse $x^2+3y^2=3$, then the maximum perpendicular distance of P from the line $x-y=10$ is -
\begin{qoptions}
\opt{A}{$3\sqrt{2}$}
\opt{B}{$4\sqrt{2}$}
\opt{C}{$6\sqrt{2}$}
\opt{D}{$5\sqrt{2}$}
\end{qoptions}
\qn{4} Two parabolas $x=cy^2+ay+b$ and $x=cy-y^2$ touch each other at the point (0,1), then the value of $2a+b-c$ is-
\begin{qoptions}
\opt{A}{$-1$}
\opt{B}{$-7$}
\opt{C}{$-5$}
\opt{D}{1}
\end{qoptions}
\qn{5} The length of common chord of the circles $x^2+y^2+2x+4y-20=0$ and $x^2+y^2+6x-8y+10=0$ is-
\begin{qoptions}
\opt{A}{$5\sqrt{\dfrac{3}{2}}$}
\opt{B}{$2\sqrt{\dfrac{3}{2}}$}
\opt{C}{5}
\opt{D}{$\dfrac{5\sqrt{5}}{2}$}
\end{qoptions}
\qn{6} A ray of light travelling along the line $x=2$ strikes a parabolic mirror with equation $y^2-2y-4x+5=0$ and gets reflected from its surface then equation of reflected ray may be -
\begin{qoptions}
\opt{A}{$y=5$}
\opt{B}{$x=-1$}
\opt{C}{$y=-1$}
\opt{D}{$y=1$}
\end{qoptions}
\qn{7} If for a hyperbola the ratio of length of conjugate Axis to the length of transverse axis is 3 : 2 then the ratio of distance between the focii to the distance between the two directrices is-
\begin{qoptions}
\opt{A}{13 : 4}
\opt{B}{4 : 13}
\opt{C}{$\sqrt{13}$ : 2}
\opt{D}{2 : $\sqrt{13}$}
\end{qoptions}
\qn{8} If $\vec{a},\vec{b}$ and $\vec{c}$ are unit vectors such that $\vec{a}$ is perpendicular to $\vec{b}$ and $\vec{c}$ and $|\vec{a}+\vec{b}+\vec{c}|=1$, then angle between $\vec{b}$ and $\vec{c}$ is-
\begin{qoptions}
\opt{A}{$\pi$}
\opt{B}{$\dfrac{\pi}{2}$}
\opt{C}{0}
\opt{D}{$\dfrac{2\pi}{3}$}
\end{qoptions}
\qn{9} Minimum distance between two points P and Q on the ellipse $\dfrac{x^2}{25}+\dfrac{y^2}{4}=1$, if difference between eccentric angles of P and Q is $\dfrac{3\pi}{2}$, is-
\begin{qoptions}
\opt{A}{$2\sqrt{2}$}
\opt{B}{$2\sqrt{5}$}
\opt{C}{$\sqrt{29}$}
\opt{D}{$\sqrt{62}$}
\end{qoptions}
\qn{10} The value of `a' for which the three given planes\par
P$_1$ : $(a+1)x-y-z=a$\par
P$_2$ : $x-(a+1)y-z=a$\par
P$_3$ : $x-y-(a+1)z=a$\par
have no point in common, is-
\begin{qoptions}
\opt{A}{0 or $-$3}
\opt{B}{$-$3}
\opt{C}{0}
\opt{D}{No such value of `a' exists}
\end{qoptions}
\qn{11} A plane passing through a point (2,2,2) cuts positive semi axes at A,B and C and P($\alpha,\beta,\gamma$) is the centroid of tetrahedron OABC (where O being origin) then select the correct option -
\begin{qoptions}
\opt{A}{$\alpha+\beta+\gamma\geq\dfrac{9}{2}$}
\opt{B}{$\alpha+\beta+\gamma=1$}
\opt{C}{$\alpha\beta\gamma=2$}
\opt{D}{$\alpha\beta\gamma\leq3$}
\end{qoptions}
\qn{12} If $\vec{p}$ and $\vec{q}$ are unequal unit vectors such that $(\vec{p}-\vec{q})\cdot((2\vec{q}+\vec{p})\times(3\vec{p}-\vec{q}))=|\vec{p}+\vec{q}|$, then angle between $\vec{p}$ and $\vec{q}$ will be-
\begin{qoptions}
\opt{A}{$\dfrac{\pi}{2}$}
\opt{B}{$\dfrac{\pi}{4}$}
\opt{C}{$\pi$}
\opt{D}{0}
\end{qoptions}
\qn{13} If $\vec{p}$ and $\vec{q}$ are unit vectors such that $[\vec{p}\ \vec{q}\ \vec{p}\times\vec{q}]=\dfrac{1}{2}$, then angle between $\vec{p}$ and $\vec{q}$ is -
\begin{qoptions}
\opt{A}{$\dfrac{\pi}{6}$}
\opt{B}{$\dfrac{\pi}{4}$}
\opt{C}{$\sin^{-1}\left(\dfrac{1}{4}\right)$}
\opt{D}{$\dfrac{\pi}{2}$}
\end{qoptions}
\qn{14} If $\vec{x}$ is a unit vector such that $\vec{x}\times(\hat{i}-2\hat{j}+\hat{k})=-\hat{i}+\hat{k}$, then $\vec{x}$ is -
\begin{qoptions}
\opt{A}{$-\hat{i}$}
\opt{B}{$-\dfrac{1}{3}(2\hat{i}-\hat{j}+2\hat{k})$}
\opt{C}{$\dfrac{1}{\sqrt{3}}(\hat{i}-\hat{j}+\hat{k})$}
\opt{D}{$\dfrac{1}{3}(2\hat{i}+\hat{j}+2\hat{k})$}
\end{qoptions}
\qn{15} The distance between directrices of the conic $\sqrt{(x-5)^2+(y-4)^2}+\sqrt{(x-3)^2+(y-2)^2}=6$ is-
\begin{qoptions}
\opt{A}{$9\sqrt{2}$}
\opt{B}{$2\sqrt{2}$}
\opt{C}{4}
\opt{D}{$\dfrac{9}{\sqrt{2}}$}
\end{qoptions}
\qn{16} Let P is a point on hyperbola $x^2-y^2=4$, which is at minimum distance from (0,$-$1) then distance of P from x-axis is-
\begin{qoptions}
\opt{A}{0}
\opt{B}{$\dfrac{1}{2}$}
\opt{C}{1}
\opt{D}{$\sqrt{2}$}
\end{qoptions}
\qn{17} If image of point P($\lambda$,2$\lambda$) in the line mirror $x-y+2=0$ is ($\mu$,3$\lambda$), then choose the correct option -
\begin{qoptions}
\opt{A}{$\mu-\lambda=0$}
\opt{B}{$\lambda+\mu=1$}
\opt{C}{$\lambda-\mu=-1$}
\opt{D}{$\lambda+\mu=0$}
\end{qoptions}
\qn{18} If angles made by lines $x^2-4xy-y^2=0$ with positive direction of x-axis are $\theta_1$ and $\theta_2$ then the value of $\sec^2(\theta_1+\theta_2)+\left|\dfrac{1}{\tan\theta_1}+\dfrac{1}{\tan\theta_2}\right|$ is -
\begin{qoptions}
\opt{A}{1}
\opt{B}{3}
\opt{C}{9}
\opt{D}{7}
\end{qoptions}
\qn{19} Equation $x^3+8y^3+24xy=64$ represents -
\begin{qoptions}
\opt{A}{Three straight lines}
\opt{B}{One straight line and one circle}
\opt{C}{One straight line and one point}
\opt{D}{Two straight lines and one point}
\end{qoptions}
\qn{20} The value of $\lambda$ for which points A(2,2,1), B(1,1,1), C($-\lambda$,2,1) and D(3,0,$-$1) are coplanar, is-
\begin{qoptions}
\opt{A}{$\lambda=-2$}
\opt{B}{$\lambda=1$}
\opt{C}{$\lambda=2$}
\opt{D}{$\lambda=-1$}
\end{qoptions}
\sect{SEC-B}
\qn{1} Let P(2,3), Q(6,0) and R($\alpha,\beta$) be three points in x-y plane such that $|PR+QR|+|PR-QR|$ is minimum then the value of $(\alpha-2\beta)$ is
\qn{2} Let A(3,5), B($-$2,$-$7) and C($\alpha,\beta$) are three points such that $\angle$ACB is a right angle and the area of triangle ABC is $\dfrac{82}{3}$ square units. Then number of such points C, is-
\qn{3} The number of integral values of `$\alpha$' for which the abscissa of point of intersection of lines $y=x+9\alpha$ and $3\alpha x+2y+9=0$ is integer, is -
\qn{4} Let `C' be the locus of centre of circles which touch the lines $x+2y-5=0$ and $2x-4y+7=0$ and area enclosed by curve `C' with the line $x-y-5=0$ is $\dfrac{P^2}{2Q^2}$; where P and Q are relative primes, then P + Q is-
\qn{5} Let a, b and c are real numbers such that $a^2+b^2+c^2=1$, then the maximum value of $(4b-3c)^2+(4a-2c)^2+(3a-2b)^2$ is -
\qn{6} Distance of the point P(1,2,1) from the plane $2x+y-z=10$ measured along the line $\dfrac{x-5}{1}=\dfrac{2y-3}{2}=\dfrac{z-\frac{5}{2}}{1}$ is -
\qn{7} If lines $\dfrac{x-1}{3}=\dfrac{y-2}{-1}=\dfrac{z-\lambda}{2}$ and $\dfrac{x+1}{-2}=\dfrac{y}{3\lambda}=\dfrac{2z-7}{1}$ are coplanar then sum of value(s) of $\lambda$ is-
\qn{8} A plane passes through a fixed point A(1,$-$2,3). If locus of foot of perpendicular to it from the point B(0,1,2) is $x^2+y^2+z^2+\alpha x+\beta y+\gamma z+\delta=0$, then the value of $|\alpha+\beta+\gamma+\delta|$ is-
\qn{9} Let mirror image of point A($\alpha,\beta$) in line mirror $x+2y=3$ is point B and image of B in lines $3x-2y=5$ is C. If origin is the orthocentre of triangle ABC and P(a,b) is point inside the triangle such that triangles PAB, PBC and PCA have same area then (a + b) is-
\qn{10} Let P($\alpha,\beta$) be a point on the parabola $y^2=4x$ which is at minimum distance from the circle $x^2+y^2-4x-20y+103=0$, then $\dfrac{\alpha}{\beta^2}$ is -
\end{multicols}
\end{document}
